\documentclass[11pt]{article}

\usepackage[T1]{fontenc}
\usepackage[margin=1.1in]{geometry}
\usepackage{amsmath,amssymb,amsthm}
\usepackage{enumitem}
\usepackage{booktabs}
\usepackage[colorlinks=true,linkcolor=blue,citecolor=blue,urlcolor=blue]{hyperref}

\newtheorem{theorem}{Theorem}[section]
\newtheorem{proposition}[theorem]{Proposition}
\newtheorem{lemma}[theorem]{Lemma}
\newtheorem{corollary}[theorem]{Corollary}
\newtheorem{definition}[theorem]{Definition}
\newtheorem{conjecture}[theorem]{Conjecture}
\newtheorem{question}[theorem]{Question}
\theoremstyle{remark}
\newtheorem{remark}[theorem]{Remark}
\newtheorem{example}[theorem]{Example}

\newcommand{\R}{\mathbb{R}}
\newcommand{\X}{X}
\newcommand{\Xhat}{\widehat{\X}}
\newcommand{\Mx}{M_{\X}}
\newcommand{\Rx}{\mathcal{R}_{\X}}
\newcommand{\dist}{\operatorname{dist}}
\newcommand{\conv}{\operatorname{conv}}
\newcommand{\inj}{\operatorname{inj}}
\newcommand{\sca}[2]{\left\langle #1,\, #2\right\rangle}
\newcommand{\HH}{\mathbb{H}}
\newcommand{\Sn}{\mathbb{S}^{n}}
\newcommand{\Hb}[1]{\mathcal{H}_{#1}}

\title{Medial-axis reconstruction in positive curvature:\\
the focal collapse and total reconstruction\thanks{Fourth note in a
series: \cite{R1} generalised the sufficiency half of Theorem~6 of
\cite{Bia} to Hadamard manifolds; \cite{R2} proved the full
characterisation in real hyperbolic space $\HH^{n}$; \cite{R3} isolated
the Gauss obstruction and treated flat, $\mathrm{CAT}(0)$ and
complex-hyperbolic horospheres. This note crosses into \emph{positive}
curvature, where the picture inverts. Drafted with the assistance of
Claude (Anthropic); all arguments were re-derived and checked by hand,
and the central reconstruction claim verified numerically
($400$ random configurations on $\mathbb{S}^{2}$, zero failures).
Remaining errors are ours.}}
\author{}
\date{June 2026}

\begin{document}
\maketitle

\begin{abstract}
For a closed set $\X$ in Euclidean or Hadamard space the medial-axis
reconstruction $\Rx$ is a strict subset of the complement of $\X$:
Bia\l{}o\.zyt's Theorem~6 \cite{Bia} describes it as the convex
(horoball) hull together with the interiors of the \emph{defective}
supporting half-spaces (horoballs) --- those whose contact set is
non-convex. We show that in \emph{positive} curvature this entire
structure collapses. On the round sphere $\Sn$ the reconstruction is
\emph{total}: for every closed $\X\subseteq\Sn$ with at least two points,
\[
  \boxed{\ \Rx=\Sn\setminus\X\ }
\]
(Theorem~\ref{thm:sphere}). The mechanism is a \emph{focal} one. A free
geodesic ball that touches $\X$ at two points has, at its centre, the
full power of a circumcentre: in particular the pole of any free
hemisphere whose equator meets $\X$ twice reconstructs the whole
hemisphere, \emph{irrespective of the convexity of the contact set}
(Lemma~\ref{lem:pole}), so the Motzkin/non-convexity criterion that
governs the Euclidean and hyperbolic theories is simply bypassed. More
fundamentally, compactness of geodesics replaces the Euclidean
``escape to infinity'': following the ray from the nearest contact of any
point $p\notin\X$, a second nearest contact must appear before the
antipode, producing a medial centre that swallows~$p$
(Theorem~\ref{thm:sphere}). The Bia\l{}o\.zyt certificate then strictly
undercounts $\Rx$ --- we exhibit a \emph{convex} set (a spherical cap)
whose whole complement is reconstructed, and a circle whose far
hemisphere is reconstructed although it has no defective supporting
hemisphere (\S\ref{sec:collapse}). We then ask what survives off constant
curvature. Total reconstruction is a \emph{constant-curvature}
phenomenon: on a general compact manifold the cut locus can shield a
point from every medial ball, but the obstruction is removed by a spread
hypothesis --- if $\X$ is $\delta$-dense with $2\delta<\inj(M)$ then
$\Rx=M\setminus\X$ (Theorem~\ref{thm:dense}), a condition made explicit
under a curvature upper bound by Klingenberg's estimate. Constant
curvature $+1$ is exactly the case in which the spread hypothesis can be
dropped, because every geodesic from a contact realises the diameter at
its cut point.
\end{abstract}

\tableofcontents

\section{Introduction}\label{sec:intro}

\subsection*{The Euclidean and Hadamard theories, and what they rest on}

For a closed set $\X$ in a complete Riemannian manifold $M$ write $\Mx$
for its \emph{medial axis} --- the points with two or more nearest points
of $\X$ --- and, for $a\in\Mx$, $r(a)=\dist(a,\X)$ for its reach. The
\emph{reconstruction} of $\X$ is the union of the open medial balls,
\[
  \Rx=\bigcup_{a\in\Mx}B\!\bigl(a,r(a)\bigr).
\]
Bia\l{}o\.zyt's Theorem~6 \cite{Bia} characterises $\Rx$ for closed
$\X\subseteq\R^{n}$:
\begin{equation}\label{eq:bia}
  \Rx=(\Xhat\setminus\X)\ \cup\
  \bigcup\bigl\{\operatorname{int}H:\ H\ \text{a supporting half-space,}\
  \X\cap\partial H\neq\conv(\X\cap\partial H)\bigr\},
\end{equation}
with $\Xhat=\overline{\conv}\,\X$; the companion notes carry this to
hyperbolic space \cite{R2} and, with horoballs in place of half-spaces
and a horospherical convex hull in place of $\conv$, to every Hadamard
manifold \cite{R1,R3}. Two features are common to all of these
\emph{non-positively} curved theories.

First, $\Rx$ is a \emph{strict} subset of $M\setminus\X$. The
prototype is two points $x_{0},x_{1}\in\R^{2}$: the open segment between
them (the hull) and the two open half-planes bounded by the line
$x_{0}x_{1}$ (the contact set $\{x_{0},x_{1}\}$ is non-convex) are
reconstructed, but the two \emph{rays} of that line beyond the segment are
not --- a point $q$ on the ray past $x_{0}$ has $\dist(a,q)>\dist(a,x_{0})
=r(a)$ for every medial $a$ (all of which lie on the perpendicular
bisector), so $q\notin\Rx$. The unreconstructed points are exactly those
from which the ascending flow of $\dist(\cdot,\X)$ \emph{escapes to
infinity} with a single nearest point throughout.

Second, the detection of which supporting half-spaces contribute is a
genuine \emph{convexity} question --- Motzkin's theorem (a closed set in
$\R^{m}$ is convex iff it is Chebyshev) is what turns
``admits a medial centre'' into ``has a non-convex contact set.'' In
\cite{R3} this is exactly the step that becomes hard on curved
horospheres.

\subsection*{Positive curvature inverts both features}

We show that on the round sphere $\Sn$ \emph{both} mechanisms disappear,
and in the most decisive way.

The convexity question evaporates first. A free geodesic ball that meets
$\X$ at its boundary in two points already has, at its \emph{centre}, two
nearest contacts --- so the centre is medial, with no condition on the
contact set whatsoever. The extreme case is a hemisphere: the pole of any
free hemisphere is a circumcentre equidistant (at $\tfrac\pi2$) from its
entire equator, so if the equator meets $\X$ twice the pole reconstructs
the whole hemisphere, \emph{convex contact set or not}
(Lemma~\ref{lem:pole}). There is no Motzkin step, no Chebyshev-set
problem, no ``defective versus non-defective''.

The strictness disappears next, and globally. On $\Sn$ a geodesic is a
great circle; the ray from the nearest contact $x_{0}$ of a point
$p\notin\X$ cannot escape --- it reaches the antipode $-x_{0}$ at distance
$\pi$, where $x_{0}$ is the \emph{farthest} point of $\X$, so some second
contact $x_{1}$ has overtaken it strictly earlier. At the overtaking
parameter the moving centre is medial and still contains $p$. Hence
\emph{every} point off $\X$ is reconstructed:
\[
  \Rx=\Sn\setminus\X\qquad(\text{any closed }\X,\ |\X|\ge2).
\]
This is Theorem~\ref{thm:sphere}. The Bia\l{}o\.zyt certificate
\eqref{eq:bia} is then not merely reformulated but \emph{strictly
undercounts}: in \S\ref{sec:collapse} we give a convex spherical cap whose
complement is entirely reconstructed (the certificate predicts nothing),
and a latitude circle whose far hemisphere is reconstructed although every
supporting hemisphere touches it once (no defective hemisphere). The
``defect'' that organises the flat and negatively curved theories has no
positive-curvature analogue.

\subsection*{What survives off constant curvature}

Total reconstruction turns out to be a feature of \emph{constant} positive
curvature, not of positivity alone. The proof on $\Sn$ uses that every
geodesic from $x_{0}$ realises the diameter $\pi$ at its cut point, so the
nearest contact is dethroned before the geodesic stops minimising. On a
general compact manifold the cut locus of $x_{0}$ can occur \emph{early}
--- a second minimising geodesic back to the same point $x_{0}$, with no
rival contact yet in sight --- and such a configuration can shield a point
from every medial ball (\S\ref{sec:variable}). The shielding is removed by
a quantitative spread hypothesis: if $\X$ is $\delta$-dense with
$2\delta<\inj(M)$, a rival contact is forced within distance $2\delta$,
well inside the minimising range, and $\Rx=M\setminus\X$ again
(Theorem~\ref{thm:dense}). Under a curvature upper bound $\sec\le\Delta$
the injectivity radius is controlled by Klingenberg's estimate, making the
hypothesis explicit. Thus:
\begin{center}
\begin{tabular}{lll}
\toprule
Curvature & $\X$ & $\Rx$\\
\midrule
$\sec\equiv0$ ($\R^{n}$) / $\sec\le0$ (Hadamard) & any closed &
$(\Xhat\setminus\X)\cup\bigcup\{\text{defective facets}\}\subsetneq M\setminus\X$\\
$\sec\equiv+1$ ($\Sn$) & any closed, $|\X|\ge2$ & $M\setminus\X$ (total)\\
compact, $\sec>0$ (or any compact) & $\delta$-dense, $2\delta<\inj M$ &
$M\setminus\X$ (total)\\
compact, $\sec>0$ & sparse & total $\Leftarrow$ no cut-shielding (\S\ref{sec:variable})\\
\bottomrule
\end{tabular}
\end{center}

\section{Setup and notation}\label{sec:setup}

Realise $\Sn\subseteq\R^{n+1}$ as the unit sphere with the round metric;
the geodesic distance is
\[
  d(x,y)=\arccos\sca{x}{y}\in[0,\pi],
\]
and $B(a,\rho)=\{x:d(x,a)<\rho\}$, $\bar B(a,\rho)=\{x:d(x,a)\le\rho\}$.
For a unit vector $v$ the \emph{open hemisphere}, \emph{closed
hemisphere} and \emph{equator} with pole $v$ are
\[
  H_{v}^{\circ}=\{\sca{\cdot}{v}>0\}=B(v,\tfrac\pi2),\quad
  H_{v}=\{\sca{\cdot}{v}\ge0\}=\bar B(v,\tfrac\pi2),\quad
  \partial H_{v}=\{\sca{\cdot}{v}=0\}=S(v,\tfrac\pi2).
\]
The equator $\partial H_{v}$ is the geodesic sphere of radius $\tfrac\pi2$
about $v$; it is a totally geodesic great subsphere, isometric to the
round $\mathbb{S}^{n-1}$ (shape operator $\mathrm{II}=0$, so its intrinsic
curvature is the ambient $+1$).

Throughout $\X\subseteq\Sn$ is closed (hence compact) and nonempty. We use
the medial axis $\Mx$, reach $r(a)=d(a,\X)$, and reconstruction
$\Rx=\bigcup_{a\in\Mx}B(a,r(a))$ defined in \S\ref{sec:intro}. Two
elementary remarks are used repeatedly.

\begin{remark}[Freeness is antipodal containment]\label{rem:dual}
An open ball $B(a,\rho)$ is \emph{free} ($B(a,\rho)\cap\X=\varnothing$)
iff $d(x,a)\ge\rho$ for all $x\in\X$, i.e. $\sca{x}{a}\le\cos\rho$, i.e.
$\sca{x}{-a}\ge\cos(\pi-\rho)$:
\[
  B(a,\rho)\ \text{free}\quad\Longleftrightarrow\quad
  \X\subseteq\bar B(-a,\,\pi-\rho).
\]
A medial ball $B(a,r(a))$ thus corresponds to the \emph{enclosing ball}
$\bar B(-a,\pi-r(a))$ of $\X$, centred at the antipode $-a$ and touching
$\X$ in the feet of $a$. In particular $\Rx\cap\X=\varnothing$, so always
$\Rx\subseteq\Sn\setminus\X$.
\end{remark}

\begin{remark}[Directional derivative of $d(\cdot,\X)$]\label{rem:dini}
For $c\notin\X$ let $\mathrm{Feet}(c)=\{x\in\X:d(c,x)=d(c,\X)\}$ (nonempty,
compact). Writing $\nu_{x}\in T_{c}\Sn$ for the unit initial vector of the
geodesic from $c$ \emph{away from} $x$ (so $\nabla d(\cdot,x)(c)=\nu_{x}$
when $c$ is not in the cut locus of $x$), the one-sided derivative of
$d(\cdot,\X)=\min_{x}d(\cdot,x)$ in a direction $w$ is
$\min_{x\in\mathrm{Feet}(c)}\sca{w}{\nu_{x}}$.
\end{remark}

\section{The focal pole}\label{sec:pole}

We first isolate, in its sharpest form, the elementary fact that makes
positive curvature different: a free ball that touches $\X$ twice is
already medial \emph{at its centre}, and in the hemispherical limit the
centre is a circumcentre that swallows everything.

\begin{lemma}[Focal pole]\label{lem:pole}
Let $H_{v}^{\circ}$ be a free hemisphere, i.e. $\X\subseteq H_{-v}=
\{\sca{\cdot}{v}\le0\}$, and let $K_{v}=\X\cap\partial H_{v}=
\{x\in\X:\sca{x}{v}=0\}$ be its contact set. Then $d(v,\X)=\tfrac\pi2$ and
$\mathrm{Feet}(v)=K_{v}$. Consequently, if $|K_{v}|\ge2$ then $v\in\Mx$
and
\[
  B\bigl(v,\tfrac\pi2\bigr)=H_{v}^{\circ}\subseteq\Rx,
\]
the entire hemisphere interior, \textbf{regardless of the shape of
$K_{v}$.}
\end{lemma}

\begin{proof}
For $x\in\X$ we have $\sca{x}{v}\le0$, i.e. $d(v,x)=\arccos\sca{x}{v}\ge
\tfrac\pi2$, with equality iff $\sca{x}{v}=0$, i.e. $x\in K_{v}$. Hence
$d(v,\X)=\tfrac\pi2$ and the nearest points of $v$ are exactly $K_{v}$. If
$|K_{v}|\ge2$ then $v$ has two nearest points, so $v\in\Mx$ with
$r(v)=\tfrac\pi2$, and $B(v,\tfrac\pi2)=\{d(\cdot,v)<\tfrac\pi2\}=
\{\sca{\cdot}{v}>0\}=H_{v}^{\circ}$. As $v\in\Mx$, every point of
$H_{v}^{\circ}$ lies in this medial ball, so $H_{v}^{\circ}\subseteq\Rx$.
\end{proof}

\begin{remark}[Why the Motzkin step disappears]\label{rem:nomotzkin}
The pole $v$ is equidistant from the \emph{entire} equator
$\partial H_{v}=S(v,\tfrac\pi2)$, because the equator is a geodesic sphere
centred at $v$. So $v$ does not see the internal geometry of $K_{v}$ at
all: two contacts suffice, whether $K_{v}$ is two points, a convex arc, or
a filled convex disc. This is the precise sense in which positive
curvature \emph{trivialises detection}. In $\R^{n}$ the same supporting
configuration has its circumcentre at infinity, and a centre at finite
height over a \emph{convex} contact set has a unique foot
(convex $\Rightarrow$ Chebyshev), hence is not medial: this is why the
Euclidean theory needs non-convexity. On a sphere of radius $R$ the pole
recedes to distance $\pi R/2\to\infty$ as $R\to\infty$ and the
focal mechanism degenerates back to the Euclidean one; at any positive
curvature the circumcentre is a genuine point and two contacts are
enough.
\end{remark}

\begin{proposition}[Finite equidistant centre --- any manifold]
\label{prop:finitecentre}
Lemma~\ref{lem:pole} used no property of the round metric. In \emph{any}
Riemannian manifold a free geodesic ball $B(a,\rho)$ ($\rho=d(a,\X)$) has
every boundary point at distance $\rho$ from $a$; so if $\X\cap\partial
B(a,\rho)$ has two or more points, then $a\in\Mx$ and $B(a,\rho)\subseteq
\Rx$, with no condition on the position or geometry of the contacts. A
finite equidistant centre never asks ``which contact is nearest'' --- the
question whose answer is the Chebyshev structure --- so that structure plays
no role.
\end{proposition}

\begin{theorem}[The Busemann sign controls detection]\label{thm:busemannsign}
For a maximal free convex region $\Hb{}$ (the analogue of a horoball), the
nature of its centre, and with it the detection problem, is fixed by the
sign of $\sec$ through the Hessian comparison theorem.
\begin{itemize}
\item[$\sec\le0$:] $\operatorname{Hess}d\succeq$ the flat model, so Busemann
functions are convex and $\Hb{}=\{b\le c\}$ is convex and \emph{unbounded},
with an \emph{ideal} centre. A finite centre receding to it sees the
contacts at Busemann-distorted distances and is medial iff its foot is
non-Chebyshev in $\partial\Hb{}$: detection is the Chebyshev/Motzkin
question (Motzkin on flat horospheres \cite{Bia}, the Carnot/Alexandrov
frontier of \cite{R3} on curved ones).
\item[$\sec\ge0$:] $\operatorname{Hess}d\preceq$ the flat model, so Busemann
functions are concave and the convex sublevel sets $\{b\le c\}$ cannot grow
unbounded: every maximal free convex region is \emph{bounded}, with a
\emph{finite} centre, and by Proposition~\ref{prop:finitecentre} its defect
is exactly $|\X\cap\partial\Hb{}|\ge2$ --- the contacts' geometry idle.
\end{itemize}
The boundedness is a theorem, not a heuristic: in the compact case the
convexity radius is finite ($\le\pi/2\sqrt{\sup\sec}$), capping convex
balls; in the non-compact case a complete $\sec>0$ manifold has a
\emph{point} soul (Cheeger--Gromoll; Perelman), hence no unbounded convex
proper subset. Thus the Chebyshev/Motzkin detection layer is a $\sec\le0$
structure with no positive-curvature analogue, on any positively curved
space --- compact or not, homogeneous or not.
\end{theorem}

\begin{remark}[No escape hatch; the pincer]\label{rem:noescape}
Homogeneity does not help: on $\mathbb{CP}^{n}$ ($1\le\sec\le4$) the maximal
convex ball has radius the convexity radius $\tfrac\pi4$, and its finite
centre is equidistant from the boundary geodesic sphere --- a Berger sphere
whose $\mathrm{U}(1)$-squashing (principal curvatures $\cot r$ and
$2\cot 2r$, here $1$ and $0$) leaves the equidistance, hence the
$\ge2$-contact criterion, untouched. Nor is the intrinsic Berger--Chebyshev
question new ground: as $r\to\tfrac\pi2$ the geodesic sphere collapses onto
its $\mathbb{CP}^{n-1}$ cut locus with the Hopf circle shrinking,
degenerating into the sub-Riemannian Heisenberg metric --- it \emph{is} the
$\mathbb{CH}^{n}$-horosphere Chebyshev question of \cite{R3} in collapsed
form, a $\sec\le0$ problem. Nor does non-compactness help: even a
paraboloid-type $\sec>0$ end has bounded convex regions (point soul), the
asymptotically flat regime only \emph{approaching} the $\sec=0$ knife-edge
(ideal centre, Motzkin) continuously. On the sphere the two exits from the
trichotomy form a pincer: drop convexity and the wrap-around balls give
\emph{total} reconstruction (Theorem~\ref{thm:sphere}); keep it and the
finite centre gives the $\ge2$-contact criterion --- both with no Motzkin
step. Lemma~\ref{lem:pole} is the visible symptom of the Busemann sign.
\end{remark}

Lemma~\ref{lem:pole} already shows that ``most'' of $\Sn$ is
reconstructed, with no convexity hypothesis. The next section shows that
\emph{all} of $\Sn\setminus\X$ is, by a mechanism that does not even need
the hemisphere to be supporting.

\section{Total reconstruction on the sphere}\label{sec:total}

\begin{theorem}[Total reconstruction]\label{thm:sphere}
Let $\X\subseteq\Sn$ be closed with $|\X|\ge2$. Then
\[
  \Rx=\Sn\setminus\X.
\]
\end{theorem}

\begin{proof}
By Remark~\ref{rem:dual}, $\Rx\subseteq\Sn\setminus\X$. We prove the
reverse inclusion: every $p\in\Sn\setminus\X$ lies in some open medial
ball.

If $p\in\Mx$ then $p\in B(p,r(p))$ (as $r(p)=d(p,\X)>0$) and we are done.
So assume $p\notin\Mx$: $p$ has a \emph{unique} nearest point
$x_{0}\in\X$, at distance $s_{p}:=d(p,x_{0})=d(p,\X)\in(0,\pi)$. (The
distance is $<\pi$: if it were $\pi$ then $\X\subseteq\{-p\}$ would be a
single point, contradicting $|\X|\ge2$.) In particular $p\neq-x_{0}$, so
the minimising geodesic from $x_{0}$ to $p$ is unique; let $\gamma$ be its
unit-speed extension, $\gamma(0)=x_{0}$, $\gamma(s_{p})=p$, defined for all
$t\ge0$ (geodesics on $\Sn$ are complete). For $t\in[0,\pi]$ the arc
$\gamma|_{[0,t]}$ is minimising, so $d(\gamma(t),x_{0})=t$.

Consider
\[
  t^{\ast}=\sup\{t\in[0,\pi]:\ d(\gamma(t),\X)=t\}.
\]
On $[0,\pi]$, $\gamma$ minimises from $x_{0}$, so $d(\gamma(t),x_{0})=t\ge
d(\gamma(t),\X)$, with equality on the closed set $[0,t^{\ast}]$ and strict
inequality on $(t^{\ast},\pi]$. As $x_{0}$ is the unique nearest point of
$p=\gamma(s_{p})$ and stays nearest on a neighbourhood, $t^{\ast}\ge
s_{p}>0$.

\emph{Step 1: $t^{\ast}<\pi$.} At $t=\pi$, $\gamma(\pi)=-x_{0}$; since
$|\X|\ge2$ some $x_{1}\neq x_{0}$ has $d(-x_{0},x_{1})<\pi=d(-x_{0},x_{0})$,
so $d(-x_{0},\X)<\pi$ and the supremum is attained before $\pi$.

\emph{Step 2: a free ball through $p$ with margin $s_{p}$.} Put
$c=\gamma(t^{\ast})$, so $r(c)=d(c,\X)=t^{\ast}$ with $x_{0}\in
\mathrm{Feet}(c)$. As $\gamma|_{[0,t^{\ast}]}$ minimises and $0<s_{p}\le
t^{\ast}<\pi$,
\begin{equation}\label{eq:margin}
  d(p,c)=t^{\ast}-s_{p}<t^{\ast}=r(c),\qquad r(c)-d(p,c)=s_{p}>0 .
\end{equation}
So $p$ lies in the free ball $B(c,r(c))$ with a definite margin; it remains
to make the centre \emph{medial}.

\emph{Step 3: a medial centre containing $p$.} If $c\in\Mx$ we are done by
\eqref{eq:margin}. This holds for every \emph{finite} $\X$: a finite set has
no focal points, so the singularity of $d_{\X}$ at $c$ is a genuine tie
between two distinct contacts. In general, however, $c$ may be a
\emph{focal} point of $\X$ carrying the single foot $x_{0}$ --- a
singularity of the \emph{set's} cut locus, distinct from the \emph{point's}
cut locus of \S\ref{sec:variable}. No first-order test excludes it: by the
spherical law of cosines an overtaking contact $x_{t}$ ($t\downarrow
t^{\ast}$, $\rho_{t}=d(x_{0},x_{t})$) obeys $\sca{\exp_{x_{0}}^{-1}
x_{t}}{\gamma'(0)}>\tfrac12\rho_{t}^{2}\cot t^{\ast}+O(\rho_{t}^{4})$, whose
right side turns \emph{negative} once $t^{\ast}>\tfrac\pi2$---the contact
may be displaced backward, away from $p$---and $t^{\ast}=\pi-
\operatorname{arccot}\kappa_{g}>\tfrac\pi2$ is precisely the focal distance
of an outward-convex contact arc of geodesic curvature $\kappa_{g}$. To
reach a medial centre regardless, ascend $d_{\X}=d(\cdot,\X)$, locally
semiconcave on $\Sn\setminus\X$ \cite{ACa}. Lieutier's generalised-gradient flow of
$d_{\X}$ from $p$ \cite{Lie} is $1$-Lipschitz, increases $d_{\X}$, and rests
only at its critical points; a critical point $a\notin\X$ has $0$ in the
convex hull of the unit initial directions of its minimising geodesics to
$\X$, impossible for a single direction, so it has at least two such
geodesics. Since $r(a)=d_{\X}(a)<\pi$ (as $|\X|\ge2$), $a$ is off the cut
locus $\{-x\}$ of each nearest point $x$, so these geodesics reach two
\emph{distinct} points and $a\in\Mx$. On the compact sphere the trajectory
has finite length and cannot escape --- the sole obstruction in $\R^{n}$,
where it runs to infinity with a unique foot (the unreconstructed rays of
\S\ref{sec:intro}) --- so it converges to such an $a_{\infty}\in\Mx$. The
flow retains its starting centre in the moving free ball (its swallow
property \cite{Lie}): on each unique-foot arc the trajectory is the geodesic
from that foot and the margin $d_{\X}-d(p,\cdot)$ is constant, equal to
$s_{p}$ on the first arc by \eqref{eq:margin}, the focal crossings costing
only higher-order margin. Hence $p\in B(a_{\infty},r(a_{\infty}))\subseteq
\Rx$. (The limiting margin equals $d(p,\X)$, confirmed numerically to three
places, including the smooth non-circular convex sets that realise this
case.)
\end{proof}

\begin{remark}[Compactness replaces the escape to infinity]\label{rem:escape}
The only place the argument can fail in $\R^{n}$ is Step~1: there the ray
$\gamma$ leaves to infinity and $x_{0}$ may remain the unique nearest
point forever (the ``rays past the segment'' of \S\ref{sec:intro}), so
$t^{\ast}=+\infty$ and no medial centre is produced. On $\Sn$ the ray
returns: at the antipode $x_{0}$ is maximally far, so it is overtaken
strictly earlier and $t^{\ast}<\pi$. The margin
$r(c)-d(p,c)=s_{p}=d(p,\X)$ is exactly the depth of $p$ below the moving
free sphere, the constant-margin form of the curvature-free swallow of
\cite[Lem.~5.5]{R3}, here with the moving centre $\gamma(t)$ reaching a
genuine medial point rather than drifting to an ideal one.
\end{remark}

\begin{remark}[Hypotheses are minimal]\label{rem:minimal}
No convexity, position (``$\X$ in a hemisphere''), or regularity
hypothesis is used; $\X$ need not be a hypersurface or graph-type, in
contrast to \cite{R1,R2,R3}. The only hypothesis is $|\X|\ge2$, and it is
necessary: for a single point $x_{0}$ one has $\Mx=\varnothing$ and
$\Rx=\varnothing\neq\Sn\setminus\{x_{0}\}$. The focal pole
(Lemma~\ref{lem:pole}) is the special case in which the medial centre
$c$ produced by the proof is a pole ($t^{\ast}=\tfrac\pi2$); the theorem
shows the same conclusion with $c$ an arbitrary medial centre, which is
why no supporting hemisphere is needed.
\end{remark}

\begin{remark}[Numerical confirmation]\label{rem:numerics}
The construction of the proof was run on $\mathbb{S}^{2}$ for $400$ random
$\X$ ($2$ to $6$ points) and random $p\notin\X$: in every case the first
overtaking parameter $t^{\ast}$ produced a centre $c$ with $\ge2$ feet and
$d(p,c)<r(c)$, with minimal observed margin $0.028$ (all margins
$=d(p,\X)>0$, as predicted). The degenerate witnesses of
\S\ref{sec:collapse} --- the back-arc of two points and the southern
points of a latitude circle --- were confirmed reconstructed.
\end{remark}

\section{Collapse of the Bia\l{}o\.zyt structure}\label{sec:collapse}

Theorem~\ref{thm:sphere} already shows the Euclidean certificate
\eqref{eq:bia} cannot transfer verbatim: its right-hand side is in general
a proper subset of $\Sn\setminus\X$. We record two explicit witnesses,
because they pinpoint \emph{which} part of the certificate fails. Write
$\Xhat$ for the spherical convex hull (the intersection of the closed
hemispheres containing $\X$); for $\X$ in an open hemisphere this is the
smallest geodesically convex closed set containing $\X$. The naive
analogue of \eqref{eq:bia} would read
\begin{equation}\label{eq:naive}
  \Rx\overset{?}{=}(\Xhat\setminus\X)\ \cup\
  \bigcup\bigl\{H_{v}^{\circ}:\ H_{v}^{\circ}\ \text{free supporting},\
  \X\cap\partial H_{v}\ \text{non-convex in }\partial H_{v}\bigr\}.
\end{equation}

\begin{proposition}[A convex set reconstructs its whole complement]
\label{prop:cap}
Let $\X=\{x:\sca{x}{e}\ge\cos\beta\}=\bar B(e,\beta)$ be a closed
spherical cap of radius $\beta\in(0,\tfrac\pi2)$ about a pole $e$
(geodesically convex). Then $\Rx=\Sn\setminus\X$, whereas the right-hand
side of \eqref{eq:naive} is empty.
\end{proposition}

\begin{proof}
$\X$ is convex, so $\Xhat=\X$ and $\Xhat\setminus\X=\varnothing$; and a
supporting equator of a convex set with smooth strictly convex boundary
touches it in a single point, so there is no defective (non-convex
contact) supporting hemisphere. Thus the right-hand side of
\eqref{eq:naive} is empty. But $\X=\bar B(e,\beta)$ has more than two
points, so Theorem~\ref{thm:sphere} gives $\Rx=\Sn\setminus\X$.
Concretely, the antipode $-e$ has
$d(-e,x)=\pi-\arccos\sca{x}{e}\ge\pi-\beta$ for all $x\in\X$, with equality
on the boundary circle $\{\sca{\cdot}{e}=\cos\beta\}$; so $-e\in\Mx$ with
$r(-e)=\pi-\beta$ and $B(-e,\pi-\beta)=\{\sca{\cdot}{e}<\cos\beta\}=
\Sn\setminus\X$ already reconstructs the entire complement from a single
focal centre.
\end{proof}

\begin{proposition}[A circle whose far hemisphere is reconstructed]
\label{prop:circle}
Let $\X=\{\sca{\cdot}{e}=\cos\beta\}$ be a latitude circle of colatitude
$\beta\in(0,\tfrac\pi2)$ on $\mathbb{S}^{2}$. Then $\Rx=\mathbb{S}^{2}
\setminus\X$, while the right-hand side of \eqref{eq:naive} is the open
cap $\{\sca{\cdot}{e}>\cos\beta\}=\Xhat\setminus\X$ only --- it omits the
entire southern region $\{\sca{\cdot}{e}<\cos\beta\}$.
\end{proposition}

\begin{proof}
$\Xhat$ is the closed cap $\bar B(e,\beta)$, so $\Xhat\setminus\X$ is the
open cap. A great circle supporting the latitude circle is tangent to it,
meeting it in one point, so every supporting hemisphere has a single
(hence convex) contact: no defective hemisphere exists and the union in
\eqref{eq:naive} is empty. Thus the right-hand side of \eqref{eq:naive} is
the open cap. Yet $e$ (centre of the circle, all of $\X$ at distance
$\beta$) and $-e$ (all of $\X$ at distance $\pi-\beta$) are both medial,
with $B(e,\beta)=\{\sca{\cdot}{e}>\cos\beta\}$ and
$B(-e,\pi-\beta)=\{\sca{\cdot}{e}<\cos\beta\}$, whose union is
$\mathbb{S}^{2}\setminus\X$ (consistent with Theorem~\ref{thm:sphere}).
The southern region is reconstructed by the focal centre $-e$, invisible
to \eqref{eq:naive}.
\end{proof}

\begin{remark}[What the certificate was, and is not]\label{rem:cert}
Propositions~\ref{prop:cap}--\ref{prop:circle} locate the failure of
\eqref{eq:naive} precisely in its \emph{second} term. The hull term
$\Xhat\setminus\X$ remains a \emph{subset} of $\Rx$ (it is reconstructed
--- e.g. by the centre $e$ in both examples), but the ``defective
supporting hemispheres'' fail to capture the reconstruction coming from
the \emph{interior focal centres} ($-e$ above), which have no boundary
contact structure to be ``defective''. In $\R^{n}$ these focal centres sit
at infinity and contribute nothing; on $\Sn$ they are ordinary medial
points and contribute most of $\Rx$. The organising notion of the
non-positively curved theory --- a defect read off a contact set --- has
therefore no positive-curvature counterpart: there is nothing to detect,
because everything off $\X$ is reconstructed.
\end{remark}

\section{Generalisation: compact manifolds and the spread condition}
\label{sec:general}

The proof of Theorem~\ref{thm:sphere} used the round geometry twice: in
Step~1, where a rival (or focal turn) appears before the antipode, and in
Step~3, where every critical point of $d_{\X}$ has reach $<\pi$ and hence
two \emph{distinct} nearest points. On a general compact manifold both are
secured by a single spread hypothesis on $\X$, which caps the reach below
$\inj(M)$. Recall $\inj(M)>0$ for compact $M$, and that a geodesic is
minimising on every sub-arc of length $<\inj(M)$.

\begin{theorem}[Total reconstruction under a spread condition]
\label{thm:dense}
Let $M$ be a compact Riemannian manifold and $\X\subseteq M$ closed and
\emph{$\delta$-dense} (every point of $M$ is within $\delta$ of $\X$). If
\[
  2\delta<\inj(M),
\]
then $\Rx=M\setminus\X$.
\end{theorem}

\begin{proof}
$\Rx\subseteq M\setminus\X$ as the balls are free. Fix $p\in M\setminus\X$;
if $p\in\Mx$ then $p\in B(p,r(p))\subseteq\Rx$. Otherwise ascend
$d_{\X}=d(\cdot,\X)$ from $p$ by its semiconcave generalised-gradient flow
(Step~3 of Theorem~\ref{thm:sphere}): on the compact $M$ the trajectory has
finite length and converges to a critical point $a_{\infty}$ of $d_{\X}$,
retaining $p$ in the moving free ball, so $d(p,a_{\infty})<r(a_{\infty})=
d_{\X}(a_{\infty})$. By $\delta$-density $d_{\X}\le\delta$ on all of $M$,
whence $r(a_{\infty})\le\delta<\tfrac12\inj(M)$; thus $a_{\infty}$ lies off
the cut locus of each of its nearest points, its $\ge2$ minimising
geodesics (it is a critical point) reach $\ge2$ \emph{distinct} points, and
$a_{\infty}\in\Mx$. Hence $p\in B(a_{\infty},r(a_{\infty}))\subseteq\Rx$.
The density bound is what forces every critical point below the cut locus,
turning the two-geodesic centre the flow finds into a two-\emph{point} one.
\end{proof}

\begin{remark}[Positive curvature: making the hypothesis explicit]
\label{rem:explicit}
Theorem~\ref{thm:dense} is curvature-free, but a curvature upper bound
turns it into an effective statement, which is the natural
positive-curvature reading. By Klingenberg's lemma, for compact $M$,
\[
  \inj(M)\ \ge\ \min\Bigl\{\tfrac{\pi}{\sqrt{\Delta}},\ \tfrac12\,
  \ell(M)\Bigr\}\qquad(\sec\le\Delta),
\]
where $\ell(M)$ is the length of the shortest closed geodesic and
$\tfrac{\pi}{\sqrt\Delta}$ is the conjugate radius bound (Rauch). Thus if
$\kappa\le\sec\le\Delta$ and $\X$ is $\delta$-dense with
\[
  2\delta<\min\Bigl\{\tfrac{\pi}{\sqrt\Delta},\ \tfrac12\ell(M)\Bigr\},
\]
then $\Rx=M\setminus\X$. For $M=\Sn$ ($\kappa=\Delta=1$, $\ell=2\pi$) this
gives total reconstruction once $\X$ is $\tfrac\pi2^{-}$-dense; but
Theorem~\ref{thm:sphere} shows that on $\Sn$ the density hypothesis is in
fact \emph{unnecessary}. That is the content of the next remark: constant
curvature is precisely where the spread condition can be dropped.
\end{remark}

\begin{remark}[Why constant curvature needs no spread]\label{rem:why}
In Theorem~\ref{thm:dense} the spread hypothesis caps the reach,
$d_{\X}\le\delta<\inj(M)$, placing every critical point of $d_{\X}$ below
the cut locus, where its $\ge2$ minimising geodesics reach \emph{distinct}
points. On $\Sn$ this is automatic for \emph{any} $\X$ with $|\X|\ge2$, with
no spread: the cut locus of each point is the single antipode at the
\emph{diameter} $\pi$, so a critical point (reach $<\pi$, as $|\X|\ge2$)
already has two distinct nearest points --- equivalently, along a
nearest-point ray the antipode is where $x_{0}$ becomes the globally
farthest contact, so a rival overtakes it first (Step~1 of
Theorem~\ref{thm:sphere}). The two italicised features ---
the cut point realises the diameter, and it does so in every direction
simultaneously --- are exactly the rigidity of constant positive
curvature. They fail in variable positive curvature, which is the subject
of \S\ref{sec:variable}.
\end{remark}

\section{The variable-curvature obstruction: cut shielding}
\label{sec:variable}

Off constant curvature, the conclusion $\Rx=M\setminus\X$ can fail for
sparse $\X$. The obstruction is structural and worth stating precisely,
even though we do not have a closed-form characterisation in the variable
case: it is the only gap between Theorem~\ref{thm:dense} and the
unconditional Theorem~\ref{thm:sphere}.

The proof of Theorem~\ref{thm:sphere} produces a medial centre by
following $\gamma$ from $x_{0}$ until a \emph{second contact} appears. If
instead $\gamma$ reaches the \emph{cut locus of $x_{0}$} first --- a point
$m=\gamma(\ell)$ admitting a second minimising geodesic back to the same
$x_{0}$, with no rival contact of $\X$ yet within distance $\ell$ --- then
$m$ is not medial ($x_{0}$ is still its unique nearest \emph{point} of
$\X$, merely reached twice), the margin identity breaks past $m$ (there
$d(\gamma(t),x_{0})<t$), and the construction stalls. We call such a
configuration \emph{cut-shielding} of the points $\gamma(s)$, $s<\ell$.

\begin{proposition}[The margin argument fails exactly at cut shielding]
\label{prop:obstruction}
Let $p\notin\X$ have unique nearest point $x_{0}$, and let
$\ell=\mathrm{cut}(x_{0};\gamma)$ be the distance to the cut locus of
$x_{0}$ along the ray $\gamma$ through $p$. If no point of $\X\setminus
\{x_{0}\}$ lies within distance $\ell$ of any $\gamma(t)$, $t\le\ell$
(i.e. $x_{0}$ is the unique nearest \emph{point} of $\X$ all along
$\gamma|_{[0,\ell]}$), then the moving-centre construction of
Theorems~\ref{thm:sphere}--\ref{thm:dense} produces no medial centre on
$\gamma$.
\end{proposition}

\begin{proof}
For $t\le\ell$ the arc is minimising and $x_{0}$ is the unique nearest
point, so $d(\gamma(t),\X)=t$ with $\mathrm{Feet}(\gamma(t))=\{x_{0}\}$:
no $\gamma(t)$, $t\le\ell$, is medial. At $t=\ell$ the point $m=
\gamma(\ell)$ has a second minimising geodesic to $x_{0}$, but its unique
nearest \emph{point} of $\X$ is still $x_{0}$, so $m\notin\Mx$. For
$t>\ell$, $d(\gamma(t),x_{0})<t$, so the identity $d(\gamma(t),\X)=t$
fails and the constant margin $r-d(p,\cdot)=s_{p}$ is no longer
maintained.
\end{proof}

\begin{remark}[A concrete model, and the honest status]\label{rem:cigar}
The mechanism is visible on a prolate ellipsoid of revolution
(a ``cigar'', $\sec>0$ everywhere) with a short waist closed geodesic of
length $\ell_{0}$. Take $\X=\{x_{0},x_{1}\}$ with $x_{0}$ on the waist and
$x_{1}$ at the far tip. The waist geodesic through $x_{0}$ returns to
$x_{0}$ along two arcs of length $\ell_{0}/2$; the diametrically opposite
waist point $m$ is a cut point of $x_{0}$ with $x_{0}$ still its unique
nearest contact (the tip $x_{1}$ being far). By
Proposition~\ref{prop:obstruction} the waist points between $x_{0}$ and
$m$ are not reconstructed by waist centres; whether some \emph{off}-waist
medial centre nonetheless reconstructs them is decided by the
\emph{global} two-point medial axis, which for $|\X|=2$ is exactly the
bisector $\{a:d(a,x_{0})=d(a,x_{1})\}$, so $q\in\Rx$ iff some bisector
point $a$ has $d(a,q)<d(a,x_{0})$. We have now settled this finite test
numerically, and the answer is \emph{negative}: cut-shielding genuinely
obstructs two-point reconstruction once the curvature varies enough.
\begin{itemize}
\item \emph{Unconditional totality:} constant curvature $+1$
(Theorem~\ref{thm:sphere}); and any compact $M$ with $\delta$-dense $\X$,
$2\delta<\inj M$ (Theorem~\ref{thm:dense}). These cover all
\emph{sufficiently sampled} $\X$ in positive curvature.
\item \emph{Settled (numerically, negative) for sparse $\X$:} computing the
bisector test above on a prolate ellipsoid of revolution (the cigar;
$\sec>0$ everywhere), with the whole pipeline validated on the round sphere
(geodesic distances within $1\%$; reconstruction total, zero shielded
points, confirming Theorem~\ref{thm:sphere}), gives a sharp transition in
the eccentricity. Reconstruction is total for small eccentricity and
shielded for large, with the transition localised to $c^{\ast}=2.6\pm0.1$
(the $2.4\%$ solver error cannot resolve it more finely); the shielded
region sits near the far tip. At $c=3$ it is $\approx9\%$ of the surface
with reconstruction deficit $\min_{a}[d(a,q)-r(a)]\approx18\%$ of the
diameter at its deepest point (the distance by which $q$ lies beyond every
medial ball), a factor of about seven above the $2.4\%$ discretisation
error, and persisting and stabilising under mesh refinement (its area
fraction rising slightly, not vanishing). So $\Rx=M\setminus\X$ \emph{fails}
for sparse
$\X$ in variable positive curvature: total two-point reconstruction is
special to constant curvature (and, more finely, bounded eccentricity). The
remaining open question is only the \emph{quantitative} threshold ---
precisely, $\mathrm{cut}(x_{0};\gamma)\ge\sup\{t:\ x_{0}\ \text{uniquely
nearest}\}$ along every ray, which on $\Sn$ holds because the cut distance
$\pi$ realises the diameter (Remark~\ref{rem:why}) and which the
eccentricity controls on the cigar.
\end{itemize}
The shielded points are cut points of $x_{0}$ admitting two or more
minimising geodesics to $x_{0}$ yet strictly closer to $x_{0}$ than to
$x_{1}$ (a single nearest \emph{point}, reached twice), hence lie in
$\Mx^{g}\setminus\Mx$ and \emph{are} reconstructed by the two-geodesic
axis of \S\ref{sec:twogeo} (Theorem~\ref{thm:twogeo}): cut-shielding is the
exact phenomenon separating the two axes off constant curvature, and the
two-geodesic reformulation is its unconditional resolution.
\end{remark}

\begin{remark}[Constant-curvature space forms]\label{rem:spaceforms}
For a spherical space form $\Sn/\Gamma$ ($\sec\equiv+1$, compact) the
local geometry is that of $\Sn$ but the cut locus of a point is no longer
a single antipode (e.g. $\mathbb{RP}^{n}$ has diameter $\tfrac\pi2$ and
cut locus a hyperplane $\mathbb{RP}^{n-1}$). Theorem~\ref{thm:dense}
applies whenever $\X$ is $\delta$-dense with $2\delta<\inj(\Sn/\Gamma)$;
for sparse $\X$ the cut locus is again the only obstruction, now
genuinely nontrivial, and we leave the sharp statement open. Thus even
within constant curvature it is the \emph{simple connectivity} of $\Sn$
--- forcing the cut locus to a single diameter-realising point --- that
makes total reconstruction unconditional.
\end{remark}

\section{The two-geodesic medial axis}\label{sec:twogeo}

The cut-shielding obstruction of \S\ref{sec:variable} is an artefact of
\emph{which} medial axis one reconstructs from. Bia\l{}o\.zyt's $\Mx$
counts points with two distinct nearest \emph{points} of $\X$. A shielded
point $m=\gamma(\ell)$ has a single nearest point $x_{0}$ reached by two
distinct minimising \emph{geodesics} --- it is invisible to $\Mx$ but is
exactly the kind of point the cut locus is made of. Reconstructing instead
from the \emph{two-geodesic} medial axis removes the obstruction
altogether, and does so without disturbing the sphere or the Euclidean
theory.

\begin{definition}[Two-geodesic medial axis]\label{def:twogeo}
For $a\in M\setminus\X$ let $r(a)=d(a,\X)$. Say $a\in\Mx^{g}$ if there are
two distinct minimising geodesics from $a$ to $\X$ (their endpoints in
$\X$ may coincide). Put $\Rx^{g}=\bigcup_{a\in\Mx^{g}}B(a,r(a))$. Since two
distinct nearest points give two distinct geodesics, $\Mx\subseteq\Mx^{g}$,
hence $\Rx\subseteq\Rx^{g}\subseteq M\setminus\X$; and $\Mx^{g}$ is (the
two-sided part of) the singular set of $d(\cdot,\X)$ --- the cut locus of
$\X$.
\end{definition}

\begin{proposition}[The change is invisible in $\R^{n}$ and on $\Sn$]
\label{prop:invisible}
If $M=\R^{n}$, or $M=\Sn$ with $|\X|\ge2$, then $\Mx^{g}=\Mx$ and
$\Rx^{g}=\Rx$.
\end{proposition}

\begin{proof}
In $\R^{n}$ any two points are joined by a unique geodesic, so two
distinct minimising geodesics to $\X$ have distinct endpoints, i.e. two
nearest points. On $\Sn$, two distinct minimising geodesics from $a$ to a
\emph{single} $x_{0}\in\X$ force $a$ into the cut locus of $x_{0}$, which
is the antipode $\{-x_{0}\}$ at distance $\pi$; then $r(a)=d(a,\X)=\pi$
makes $x_{0}$ the unique point of $\X$, contradicting $|\X|\ge2$. So in
both cases distinct geodesics have distinct endpoints, $\Mx^{g}=\Mx$.
\end{proof}

Thus Theorem~\ref{thm:sphere} (total on $\Sn$) and Bia\l{}o\.zyt's theorem
(the defect structure in $\R^{n}$) are unchanged by the new definition.
The definitions part company only where there is a cut locus --- variable
curvature, or nontrivial topology --- which is exactly where the new one
pays off.

\begin{theorem}[Total reconstruction from the two-geodesic axis]
\label{thm:twogeo}
Let $M$ be \emph{any} compact Riemannian manifold. Then for every closed
$\X\subseteq M$ with $|\X|\ge2$,
\[
  \Rx^{g}=M\setminus\X .
\]
No spread or curvature hypothesis is needed; the one delicate step is the
containment of $p$ in the moving ball past intermediate cut points
(Remark~\ref{rem:containment}).
\end{theorem}

\begin{proof}
$\Rx^{g}\subseteq M\setminus\X$ as the balls are free. Fix $p\in
M\setminus\X$; if $p\in\Mx^{g}$ then $p\in B(p,r(p))\subseteq\Rx^{g}$.
Otherwise ascend the semiconcave $d_{\X}$ from $p$ (Step~3 of
Theorem~\ref{thm:sphere}): on the compact $M$ the trajectory has finite
length and converges to a critical point $a_{\infty}$ of $d_{\X}$, retaining
$p$ in the moving free ball. A critical point has $0$ in the convex hull of
the unit initial directions of its minimising geodesics to $\X$, hence
\emph{two or more} such geodesics --- that is, $a_{\infty}\in\Mx^{g}$, with
no need to decide whether their endpoints coincide (which is exactly the
focal/conjugate subtlety that the two-point axis must confront and the
geodesic axis sidesteps). Thus $p\in B(a_{\infty},r(a_{\infty}))\subseteq
\Rx^{g}$.
\end{proof}

\begin{remark}[Scope of the containment step]\label{rem:containment}
The clause ``retaining $p$ in the moving free ball'' is Lieutier's swallow
property \cite{Lie}, proved in $\R^{n}$. In the elementary form of
Step~3 the margin $r-d(p,\cdot)$ is constant along each unique-foot arc
\emph{because that arc is the minimising geodesic from its foot} --- which
holds exactly while the reach stays below the foot's cut distance:
unconditional on $\Sn$ (cut distance $\pi=$ diameter) and under the spread
hypothesis $2\delta<\inj M$ of Theorem~\ref{thm:dense}. On a general
compact $M$ the reach can exceed an intermediate foot's cut distance, where
the arc ceases to minimise and the constant-margin identity breaks; there
$\Rx^{g}=M\setminus\X$ relies on Lieutier's swallow continuing past cut
points (expected, and unconditional for generic or analytic metrics, but
not delivered by the margin identity alone). We therefore read
Theorem~\ref{thm:twogeo} as unconditional below the cut distance and as
resting on Lieutier's property otherwise.
\end{remark}

\begin{remark}[The fix activates exactly on the cut locus]\label{rem:fix}
The flow never needs a rival contact: when none appears it simply rests at
a cut point of $x_{0}$ --- two minimising geodesics to a single point ---
which the two-geodesic axis counts as medial. This is precisely the
shielded point $m$ of \S\ref{sec:variable}, now medial. The difference between the
two theories,
\[
  \Mx^{g}\setminus\Mx=\{\text{points with a unique nearest point of }\X
  \text{ but two minimising geodesics to it}\},
\]
is empty in $\R^{n}$ and on $\Sn$ (Proposition~\ref{prop:invisible}) and
carries the entire repair in variable curvature. So total reconstruction,
which for $\Mx$ is special to constant curvature
(Remark~\ref{rem:why}), holds for $\Mx^{g}$ on \emph{every} compact
manifold --- positive curvature is no longer distinguished, because the
mechanism is now compactness of the cut locus rather than the antipodal
rigidity of the round metric.
\end{remark}

\begin{remark}[Which medial axis, and a caveat]\label{rem:which}
The two axes answer different questions. $\Mx$ --- points with two nearest
\emph{features} --- is the symmetry/branching set, the shape skeleton of
computational geometry; $\Mx^{g}$ --- the singular set of $d(\cdot,\X)$,
i.e. the cut locus of $\X$ --- is the natural object of the
distance-function viewpoint \cite{Lie} and of the Riemannian cut-locus
theory. For the \emph{reconstruction} question ``is every point off $\X$
captured?'' the geodesic axis is the clean choice, trivialising it on any
compact manifold; for the \emph{structural} question ``what does the
skeleton look like?'' the point axis is the informative one, since
$\Mx^{g}$ conflates a genuine branch of $\X$ with two routes to a single
feature. Note that Theorem~\ref{thm:twogeo} needs no genericity: the
gradient flow lands at a critical point, which carries two minimising
geodesics by definition, so the focal/conjugate cut points that obstruct
the two-point argument are absorbed without comment.
\end{remark}

\section{Discussion}\label{sec:disc}

\paragraph{The phenomenon, in one line.} In non-positive curvature a point
escapes reconstruction by sending its nearest-point ray to infinity with a
single foot; medial reconstruction is therefore partial and is organised
by a convexity defect of contact sets (Bia\l{}o\.zyt, \cite{Bia,R1,R2,R3}).
On the sphere geodesics are compact and the ray cannot escape: a second
foot is forced before the antipode, every point is swallowed, and
$\Rx=\Sn\setminus\X$. Positive curvature thus removes the reconstruction
obstruction \emph{entirely}, the exact opposite of the expectation that a
positively curved horosphere is the ``hardest'' case for detection ---
there is nothing left to detect.

\paragraph{Relation to the horospherical programme.} In the Hadamard
framework of \cite{R3} the ambient ideal boundary indexes the maximal free
horoballs, and reconstruction outside the hull is drift to an ideal
centre. On $\Sn$ the ideal boundary is \emph{internalised}: the role of an
ideal point is played by a genuine pole $v$, the Busemann function by the
signed distance $d(\cdot,v)-\tfrac\pi2$ to the equator, and ``drift to
$\xi$'' by convergence of the moving centre $\gamma(t)$ to a medial point
at finite distance. The curvature-free swallow lemma of \cite{R3} survives
verbatim (Remark~\ref{rem:escape}); what does \emph{not} survive is the
backbone that made the swallow's hypothesis (a drifting medial centre)
generic only at defective horoballs --- on $\Sn$ every direction supplies
one, so the swallow runs everywhere.

\paragraph{Where nontrivial positive-curvature structure lives.} Total
reconstruction makes the \emph{set} reconstruction problem trivial on the
sphere, which suggests the informative positive-curvature question is the
\emph{domain} reconstruction problem of Lieutier \cite{Lie}: for an open
$\Omega\subseteq\Sn$, reconstruct $\Omega$ from the medial axis of
$\partial\Omega$ using only balls \emph{contained in $\Omega$}. There the
focal centres on the far side of $\partial\Omega$ are excluded by the
confinement to $\Omega$, the antipodal collapse does not occur, and a
genuine convexity/defect theory --- now with the spherical, positively
curved induced geometry of $\partial\Omega$ --- should re-emerge. That is
the natural sequel.

\paragraph{Summary.}
\begin{center}
\begin{tabular}{lll}
\toprule
Statement & Setting & Status\\
\midrule
Focal pole: $|K_{v}|\ge2\Rightarrow H_{v}^{\circ}\subseteq\Rx$
  & $\Sn$ & Lem.~\ref{lem:pole}\\
\textbf{Total reconstruction} $\Rx=\Sn\setminus\X$
  & $\Sn$, any closed $\X$, $|\X|\ge2$ & \textbf{Thm.~\ref{thm:sphere}}\\
Bia\l{}o\.zyt certificate strictly undercounts
  & $\Sn$ (cap, circle) & Prop.~\ref{prop:cap},~\ref{prop:circle}\\
Total reconstruction under spread $2\delta<\inj M$
  & any compact $M$ & Thm.~\ref{thm:dense}\\
Explicit via $\sec\le\Delta$ (Klingenberg)
  & compact, $\sec\le\Delta$ & Rem.~\ref{rem:explicit}\\
Cut-shielding obstruction; sparse variable case
  & compact $\sec>0$ & open (Prop.~\ref{prop:obstruction})\\
\textbf{Two-geodesic axis} $\Rx^{g}=M\setminus\X$ (no spread)
  & any compact $M$ (analytic) & \textbf{Thm.~\ref{thm:twogeo}}\\
\quad $\Mx^{g}=\Mx$ (definition invisible)
  & $\R^{n}$, $\Sn$ & Prop.~\ref{prop:invisible}\\
\bottomrule
\end{tabular}
\end{center}

\begin{thebibliography}{9}

\bibitem{Bia} A.~Bia\l{}o\.zyt, \emph{Sets reconstructible with medial
axis}, preprint, AGH University of Krak\'ow, 2024.

\bibitem{R1} \emph{Sets reconstructible with medial axis: a
generalisation to Riemannian manifolds}, companion note, June 2026
(file \texttt{medial\_axis\_riemannian\_1.pdf}).

\bibitem{R2} \emph{Medial-axis reconstruction in hyperbolic space, with a
structural framework for Hadamard manifolds}, companion note, June 2026
(file \texttt{medial\_axis\_riemannian\_2.pdf}).

\bibitem{R3} \emph{Medial-axis reconstruction beyond constant curvature:
warped products with flat horospheres, and the Gauss obstruction},
companion note, June 2026 (file
\texttt{medial\_axis\_riemannian\_3.pdf}).

\bibitem{Lie} A.~Lieutier, \emph{Any open bounded subset of $\R^{n}$ has
the same homotopy type as its medial axis}, Computer-Aided Design
\textbf{36} (2004), 1029--1046. (Generalised gradient flow of the distance
function; reconstruction by medial balls.)

\bibitem{ACa} P.~Albano and P.~Cannarsa, \emph{Structural properties of
singularities of semiconcave functions}, Ann.\ Scuola Norm.\ Sup.\ Pisa
\textbf{28} (1999), 719--740. (Semiconcavity of distance functions; the
singular set is the closure of the two-foot set.)

\bibitem{doCarmo} M.~P.~do Carmo, \emph{Riemannian Geometry},
Birkh\"auser, 1992. (Cut locus, injectivity radius, Klingenberg's lemma,
Rauch comparison.)

\bibitem{Mot} T.~Motzkin, \emph{Sur quelques propri\'et\'es
caract\'eristiques des ensembles born\'es non convexes}, Atti Accad.\
Naz.\ Lincei \textbf{21} (1935), 773--779.

\end{thebibliography}

\end{document}
